% =====================================================================
%  CV w LaTeXu -- jedna strona A4
%  Kompilacja:  pdflatex cv.tex   (albo: latexmk -pdf cv.tex, Overleaf)
% =====================================================================
\documentclass[10pt,a4paper]{article}

% ---------- kodowanie, jezyk, czcionka ----------
% dziala zarowno w pdfLaTeX (Overleaf), jak i w XeLaTeX/LuaLaTeX
\usepackage{iftex}
\ifPDFTeX
  \usepackage[T1]{fontenc}
  \usepackage[utf8]{inputenc}
  \usepackage{lmodern}
\else
  \usepackage{fontspec}   % domyslnie Latin Modern z pelnym Unicode
\fi
\usepackage[polish]{babel}

% ---------- uklad strony ----------
\usepackage[a4paper,top=1.6cm,bottom=1.4cm,left=1.8cm,right=1.8cm]{geometry}
\usepackage{xcolor}
\usepackage{titlesec}
\usepackage{enumitem}
\usepackage{tabularx}
\usepackage[hidelinks]{hyperref}

\definecolor{akcent}{RGB}{140,30,40}
\pagestyle{empty}
\setlength{\parindent}{0pt}

% naglowki sekcji: kolorowy tytul + cienka linia
\titleformat{\section}
  {\large\scshape\color{akcent}}{}{0pt}{}[\vspace{-0.7ex}{\color{akcent}\titlerule}]
\titlespacing*{\section}{0pt}{2.2ex}{1.2ex}

\setlist[itemize]{leftmargin=1.2em,itemsep=1pt,topsep=2pt,parsep=0pt}

% \pozycja{Tytul}{Daty}{Miejsce / technologie}
\newcommand{\pozycja}[3]{%
  \textbf{#1}\hfill{\small #2}\\
  {\itshape #3}\par}

\usepackage{graphicx}

% ---------- DANE OSOBOWE ----------
\newcommand{\imienazwisko}{Artom Busniuk}
\newcommand{\telefon}{+48 72933132}            % TODO: sprawdz numer (brakuje cyfry?)
\newcommand{\email}{kendiplay6@gmail.com}
\newcommand{\github}{github.com/LOGIN}         % TODO: wpisz swoj login
\newcommand{\miasto}{Wrocław}
% Zdjecie: wrzuc plik zdjecie.jpg do folderu -- pojawi sie automatycznie.
% Bez pliku naglowek jest wysrodkowany, bez zdjecia.
\newcommand{\zdjecie}{zdjecie.jpg}

\begin{document}

% =====================================================================
\IfFileExists{\zdjecie}{%
  \begin{minipage}[c]{0.74\linewidth}
    {\fontsize{26}{30}\selectfont\bfseries \imienazwisko}\\[2ex]
    \small
    \miasto\\[0.3ex]
    tel.\ \telefon\\[0.3ex]
    \href{mailto:\email}{\email}\\[0.3ex]
    \href{https://\github}{\github}
  \end{minipage}\hfill
  \begin{minipage}[c]{0.24\linewidth}\raggedleft
    \includegraphics[width=3cm,height=3.8cm,keepaspectratio]{\zdjecie}
  \end{minipage}
}{%
  \begin{center}
    {\fontsize{26}{30}\selectfont\bfseries \imienazwisko}\\[1.5ex]
    \small
    \miasto\quad\textbar\quad
    \telefon\quad\textbar\quad
    \href{mailto:\email}{\email}\quad\textbar\quad
    \href{https://\github}{\github}
  \end{center}
}

% =====================================================================
\section*{O mnie}
Student pierwszego roku informatyki na Uniwersytecie Wrocławskim. Programuję
w~Pythonie, poznaję programowanie funkcyjne i~algorytmikę. Szukam pierwszych
praktyk lub stażu, na którym będę mógł uczyć się od doświadczonych programistów.

% =====================================================================
\section*{Wykształcenie}
\pozycja{Informatyka, studia inżynierskie (I~stopnia)}{od października 2026}
  {Instytut Informatyki, Uniwersytet Wrocławski}

% =====================================================================
\section*{Projekty}
\pozycja{Wyszukiwarka palindromów w~tekście}{listopad 2026}{OCaml, Perl}
\begin{itemize}
  \item Program w~OCamlu znajdujący najdłuższy palindrom w~polskiej powieści
        (ok.~600~KB, UTF-8) w~czasie liniowym przy użyciu algorytmu Manachera.
  \item Jednolinijkowiec w~Perlu (17~znaków) sumujący wszystkie liczby w~pliku.
\end{itemize}
\medskip

\pozycja{CV w~\LaTeX{}u}{listopad 2026}{\LaTeX, Git, GitHub}
\begin{itemize}
  \item Niniejsze CV, prowadzone w~repozytorium Git (wyłącznie pliki źródłowe).
\end{itemize}

% =====================================================================
\section*{Umiejętności}
\begin{tabularx}{\linewidth}{@{}l X@{}}
  \textbf{Programowanie:} & Python; podstawy: OCaml, Perl \\
  \textbf{Narzędzia:}     & Git, GitHub, Linux / macOS (terminal), \LaTeX \\
  \textbf{Języki:}        & polski -- ojczysty, angielski -- B2 \\
\end{tabularx}

% =====================================================================
\vfill
{\scriptsize\color{gray}
Wyrażam zgodę na przetwarzanie moich danych osobowych dla potrzeb niezbędnych
do realizacji procesu rekrutacji zgodnie z~Rozporządzeniem Parlamentu Europejskiego
i~Rady (UE) 2016/679 z~dnia 27 kwietnia 2016~r. w~sprawie ochrony osób fizycznych
w~związku z~przetwarzaniem danych osobowych i~w~sprawie swobodnego przepływu takich
danych oraz uchylenia dyrektywy 95/46/WE (RODO).\par}

\end{document}
